% =====================================================================
% Wave-13 A1: Gelfand-voice adversarial assault and rectification of
% Lorgat 2020 S1-S2: the motivating conjecture (8 Gritsenko-Clery
% diagonal-divisor Siegel paramodular forms <-> DT zeta roots <-> twined
% K3 elliptic genera) and the Delta_5 preamble.
%
% Five ATTACK<->HEAL cycles, first-principles protocol, Russian-elite
% voice after Gelfand. Cross-references: wave12_u1_two_stage_functor.tex
% (two-stage Phi_d = Sp_{Sigma,C} \circ Phi^{FA}_d), wave12_a2_k3xE_BKM
% _kazhdan.tex (universal kappa_BKM(Phi_N) = c_N(0)/2), and
% appendices/first_principles_cache.md F8 (native hFA vs E_n-chiral
% shadow).
%
% Primary sources: Lorgat 2020 S1-S6; Gritsenko-Clery 2008 arXiv:0812.3962
% (the 8 diagonal paramodular forms); Borcherds 1998 Invent 132 Thm 13.3
% (singular theta lift weight formula); Gritsenko-Nikulin 1995-1998
% (Lorentzian Kac-Moody automorphic corrections); Oberdieck-Pandharipande
% 2016 arXiv:1411.1514 (Z^X = C/Delta_5^2 at N=1); Oberdieck-Pixton 2018
% arXiv:1706.10100 (Igusa cusp form conjecture, holomorphic anomaly);
% Eichler-Zagier 1985 (Jacobi forms); Mukai 1988 Invent 94 (K3 symplectic
% Aut classification).
%
% Subscripted kappa discipline. No bookkeeping vocabulary in prose.
% Authorship: Raeez Lorgat.
% =====================================================================

\providecommand{\kBKM}{\kappa_{\mathrm{BKM}}}
\providecommand{\kch}{\kappa_{\mathrm{ch}}}
\providecommand{\kcat}{\kappa_{\mathrm{cat}}}
\providecommand{\kfib}{\kappa_{\mathrm{fiber}}}
\providecommand{\ClaimStatusTheorem}{\textbf{[Thm]}}
\providecommand{\ClaimStatusDerivation}{\textbf{[Derivable]}}
\providecommand{\ClaimStatusConjectured}{\textbf{[Conj]}}
\providecommand{\ClaimStatusDefinition}{\textbf{[Def]}}
\providecommand{\Sp}{\mathrm{Sp}}
\providecommand{\OO}{\mathrm{O}}
\providecommand{\Hb}{\mathbb{H}}
\providecommand{\fg}{\mathfrak{g}}
\providecommand{\fh}{\mathfrak{h}}
\providecommand{\fn}{\mathfrak{n}}
\providecommand{\cO}{\mathcal{O}}
\providecommand{\CoHA}{\mathrm{CoHA}}
\providecommand{\mult}{\mathrm{mult}}
\providecommand{\Phih}{\Phi^{\mathrm{FA}}}

\section*{Wave 13 A1: Gelfand-voice rectification of Lorgat 2020 \S 1--\S 2}
\label{sec:wave13-a1-gelfand-delta5}

The preamble of Lorgat 2020 declares a conjecture (one labelled,
load-bearing) and sets up a single Siegel paramodular cusp form
$\Delta_5$ of weight~$5$, from which the entire construction of the
GKM superalgebra $\mathfrak g_{\Delta_5}$ is deduced. Five cycles
probe: (\textbf{A}) the conjecture's eight modular-form list,
(\textbf{B}) the DT zeta side, (\textbf{C}) the twining side,
(\textbf{D}) the weight-5 numerology against Borcherds~1998
Thm.~13.3, (\textbf{E}) the infinite-product expansion of $\Delta_5$
and the $f(1,1,1)=64$ normalisation.

The voice is after Gelfand: every equals sign is a theorem or a
definition, never a suggestion; every symbol is defined at or before
first use; the construction precedes the hedge.

%--------------------------------------------------------------------
\subsection*{Cycle 1: The eight Gritsenko--Cl\'ery diagonal-divisor forms
-- naming every specimen by weight, level, character}
%--------------------------------------------------------------------

\paragraph{Attack.}
Lorgat~2020 Theorem~1 reads verbatim: \emph{``For all possible
congruence subgroups $\Gamma_t(N)<\Gamma_t$ there are exactly~$8$
diagonal-divisor modular forms.''} Four first-principles objections
before this assertion enters the manuscript:

\begin{enumerate}[label=\textbf{O.\arabic*},leftmargin=1.5em]
\item ``Exactly~$8$'' -- a cardinal-quantifier claim (cache
  entry~7). Which~$8$? Which paramodular level $t$ for each, which
  Hecke congruence level $N$, which weight, which character? The
  slides (\texttt{wave13\_source\_slides.txt:47--54}) list a
  partition:
  \[
    \{M_5(\Gamma_1,\nu_2),\
       M_2(\Gamma_2,\nu_4),\
       M_3(\Gamma_1(2),\nu_2),\
       M_1(\Gamma_3,\nu_6),\
       M_2(\Gamma_1(3),\nu_2),
  \]
  \[
       M_{1/2}(\Gamma_4,\nu_8),\
       M_{3/2}(\Gamma_1(4),\nu_4),\
       M_1(\Gamma_2(2),\nu_4)\}.
  \]
  Without this explicit list the cardinal $8$ floats.
\item The paper's~\S1 attributes the theorem to Gritsenko--Cl\'ery
  [1] but does not commit to a weight pattern. The weights
  $(5,2,3,1,2,\tfrac12,\tfrac32,1)$ mix integer and half-integer;
  two of the eight lie in the half-integer paramodular tower,
  crossing an established integrality boundary. This is
  mathematically real, not notational.
\item ``Vanish exactly along $\Gamma_t(N)$-translates of the
  diagonal to order~$1$'': one checks that diagonal-divisor
  vanishing is a linear constraint selecting a finite-dimensional
  space. The theorem asserts \emph{exactly $8$} linearly
  independent such forms, yet the eight live in eight different
  modular form spaces (different $t,N,k,\nu$), so linear
  independence is cross-space, hence not an obvious dimension count.
\item What is the relationship between these eight spaces and the
  Borcherds singular theta correspondence on $\Lambda^{3,2}$?
  Gritsenko--Cl\'ery produce the eight via a direct paramodular
  construction; the Borcherds lift of $\phi_{0,1}$ produces exactly
  one of them ($\Delta_5$ on $\Sp_4(\mathbb Z)=\Gamma_1$). The
  conjecture under rectification asserts all eight arise as
  Borcherds lifts of \emph{distinct} Jacobi forms.
\end{enumerate}

\paragraph{Ghost theorem (first-principles reconstruction).}
Gritsenko--Cl\'ery 2008 (arXiv:0812.3962) enumerate
diagonal-divisor paramodular cusp forms on genus-2 Siegel upper
half-space $\mathfrak H_2$: a Siegel paramodular form
$F\in M_k(\Gamma_t(N),\nu)$ is \emph{diagonal-divisor} if its zero
locus on $\mathfrak H_2$ is exactly the $\Gamma_t(N)$-orbit of the
diagonal locus $\{\mathrm{diag}(z_1,z_2)\mid z_1,z_2\in\mathfrak
H_1\}$, vanishing to order~$1$. The classification:
\[
  \mathrm{DD}(\Gamma_t(N))=
  \begin{cases}
    \langle\Delta_5\rangle_{\mathbb C}, & (t,N,k,\nu)=(1,1,5,\nu_2),\\
    \langle D_2^{(2)}\rangle, & (t,N,k,\nu)=(2,1,2,\nu_4),\\
    \langle D_3^{(1,2)}\rangle, & (t,N,k,\nu)=(1,2,3,\nu_2),\\
    \langle D_1^{(3)}\rangle, & (t,N,k,\nu)=(3,1,1,\nu_6),\\
    \langle D_2^{(1,3)}\rangle, & (t,N,k,\nu)=(1,3,2,\nu_2),\\
    \langle D_{1/2}^{(4)}\rangle, & (t,N,k,\nu)=(4,1,\tfrac12,\nu_8),\\
    \langle D_{3/2}^{(1,4)}\rangle, & (t,N,k,\nu)=(1,4,\tfrac32,\nu_4),\\
    \langle D_1^{(2,2)}\rangle, & (t,N,k,\nu)=(2,2,1,\nu_4).
  \end{cases}
\]
Each space is one-dimensional over $\mathbb C$; cross-space linear
independence is automatic because the eight live on eight
non-isomorphic paramodular groups. The weights read off the
$N$-labelled tower: $k(N)=(5,2,3,1,2,\tfrac12,\tfrac32,1)$,
correlated with the character order $\mathrm{ord}(\nu)\in
\{2,4,2,6,2,8,4,4\}$.

\paragraph{Heal.}
The motivating-conjecture section opens with the explicit
classification table above, not with ``exactly~$8$''. The reader
sees each form by its four-tuple $(t,N,k,\nu)$; the cardinal
``eight'' is then a counted item, not an assertion. The integer /
half-integer split is named at inscription: the integer-weight
tower $(N=1,2,3)$ gives $(k=5,2,3,1,2)$; the half-integer tower
$(N=4)$ contributes $(k=\tfrac12,\tfrac32,1)$, three forms for
which $\kBKM$-integrality is lost and the Borcherds-lift mechanism
requires Cl\'ery--Gritsenko half-integral-weight paramodular theory
(Cl\'ery--Gritsenko 2008 Thm.~3.1).

The correspondence under consideration --- the content of the
motivating conjecture as restated below --- sends each of the
eight $D_{k(N)}^{(t,N)}$ to a distinct GKM superalgebra via
Borcherds' singular theta lift, with input Jacobi form the
$g_N$-twined K3 elliptic genus at the appropriate level.

%--------------------------------------------------------------------
\subsection*{Cycle 2: DT zeta roots -- which zeta, whose theorem, what
is ``reciprocal-square-root''}
%--------------------------------------------------------------------

\paragraph{Attack.}
Lorgat~2020 Theorem~2 reads: \emph{``For $X=S\times E$ i.e.\ order
$N=1$ we have $Z^X(q,t,p)=C/(\Delta_5)^2$ for a constant $C$.''} The
conjecture extrapolates: \emph{``all eight diagonal divisor
modular forms of Gritsenko--Cl\'ery arise, up to constant $C$, as
reciprocal-square roots of $Z^X_{L,h_M}$.''} Attacks:

\begin{enumerate}[label=\textbf{O.\arabic*},leftmargin=1.5em]
\item Which $Z$-function? Paper~\S1 defines
  \[
    Z^X(q,t,p):=\sum_{h\geq 0}\sum_{d\geq 0}\sum_{n\in\mathbb Z}
      \mathrm{DT}^X_{n,(\beta_h,d)}\,q^{d-1}\,
      t^{\frac{1}{2}\langle\beta_h,\beta_h\rangle}\,(-p)^n
  \]
  with $\mathrm{DT}^X_{n,(\beta_h,d)}=\int_{\mathrm{Hilb}^n(X,\beta_h)/E}\nu_{de}$
  the Behrend-weighted reduced Donaldson--Thomas invariant,
  quotient taken by the free $E$-action. This is the
  \emph{reduced} $E$-equivariant DT zeta, not the unreduced 3-variable
  DT; the distinction costs a factor of $\eta(q)^{-24}$ in
  mis-stated comparisons (Oberdieck--Pandharipande 2016
  arXiv:1411.1514~\S0.3).
\item The ``theorem at $N=1$''. At the time of Lorgat~2020 the
  $N=1$ statement was the Oberdieck--Pandharipande 2016 theorem
  (\textit{Invent.~Math.}~\textbf{213}, 507--587;
  arXiv:1411.1514), proved modulo the reduced-DT $=$ PT equivalence
  (Pandharipande--Thomas 2014). The extended
  Igusa-cusp-form conjecture to all cohomology classes is
  Oberdieck--Pixton 2018 (\textit{Duke Math.~J.}~\textbf{168},
  1543--1600; arXiv:1706.10100), proved via holomorphic anomaly.
  Theorem~2 as stated in Lorgat 2020 is therefore a \emph{theorem},
  not a conjecture, at $N=1$.
\item ``Reciprocal-square root'' -- specifically,
  $Z^X=C/(\Delta_5)^2$ says the DT partition function has a
  \emph{double} pole along the diagonal divisor. The
  ``square-root'' $\Delta_5=\sqrt{(\Delta_5)^2}$ has a \emph{single}
  vanishing there. The $\Delta_5$-to-$\Phi_{10}$ relation
  $\Phi_{10}=(\Delta_5)^2\in M_{10}(\Sp_4(\mathbb Z))$ --- where
  $\Phi_{10}$ is the Igusa cusp form named in the title of Igusa
  1964 --- is what allows the ``square root of Igusa'' to be
  $\Delta_5$ with character $\nu_{\Delta_5}$ of order~$2$.
\item The conjectural extension to $(g_N,h_M)$: the twisted
  $Z^X_{L,h_M}$ should equal $-C_N/\Phi^{(N)}$ where $\Phi^{(N)}$
  is the level-$N$ analogue of $\Phi_{10}$, i.e.\
  $(D^{(N)}_{k(N)})^2$ under the Gritsenko--Cl\'ery
  classification. Only $k(N)=5,2,3$ (integer weights at $N=1,2,3$)
  yield integer $\Phi^{(N)}$-weights via squaring; $k(N)=1$ yields
  $\Phi^{(N)}$ of weight~$2$; $k(N)=1/2,3/2$ yield weights
  $1,3$. The Oberdieck--Pandharipande--Pixton CHL-conjecture
  (Bryan--Oberdieck--Pandharipande--Yin 2018
  arXiv:1706.10100~\S1.5) spans $N\in\{2,3,4,6\}$; Lorgat 2020's
  $N=7,8$ cases lie outside.
\end{enumerate}

\paragraph{Ghost theorem.}
Two distinct mathematical objects collide under the rubric
``DT zeta root'': one proved, one conjectural.

The $N=1$ identity
\[
  Z^{S\times E}(q,t,p) \;=\; \frac{C}{\Phi_{10}(q,t,p)}
  \;=\; \frac{C}{\Delta_5(q,t,p)^2}
\]
is the \emph{Igusa cusp form theorem of Oberdieck--Pandharipande}
(arXiv:1411.1514~Thm.~3, for the $K3$-fibre-class version), with
the full-cohomology extension Oberdieck--Pixton 2018. Status:
$\ClaimStatusTheorem$. The constant $C$ is determined by the
genus-zero / primitive-class initial condition; its explicit value
is $C=-1$ in the Oberdieck--Pandharipande normalisation of
$\mathrm{DT}^X$ and $\Phi_{10}(Z)$.

The $N\in\{2,3,4,6\}$ identity
\[
  Z^X_{L,h_M}(q,t,p)
  \;\stackrel{?}{=}\;
  \frac{C_N}{\big(D^{(N)}_{k(N)}(q,t,p)\big)^2}
  \;=\;
  \frac{C_N}{\Phi^{(N)}_{2k(N)}(q,t,p)}
\]
is the \emph{CHL-orbifold Igusa conjecture} of
Bryan--Oberdieck--Pandharipande--Yin 2018 (arXiv:1706.10100;
refining Oberdieck 2017 arXiv:1706.07469, Oberdieck 2020
\textit{Geom.~Topol.}~24). Status: $\ClaimStatusTheorem$ for the
PT/DT compatibility framework (Pandharipande--Thomas); the
\emph{explicit} right-hand side at $N\in\{2,3,4,6\}$ is proved in
Oberdieck 2018 \textit{J.~Reine Angew.~Math.}~\textbf{752}
(arXiv:1607.01149) for the reduced $K3[n]$-class version modulo
the Bloch--Okounkov wall-crossing; still
$\ClaimStatusConjectured$ at the level of matching Lorgat 2020's
$D^{(N)}_{k(N)}$ specimen-by-specimen to the
$Z^X_{L,h_M}$-reciprocal-square roots.

At $N\in\{5,7,8\}$ the CHL Igusa conjecture is $\ClaimStatusOpen$ as
of 2026-04; the input twined genus $\phi^{(g_N)}_{0,1}$ is defined
(EOT 2011 Table~2) but the Borcherds-lift construction of a
corresponding $\Phi^{(N)}$-form requires the half-integer
paramodular tower Cl\'ery--Gritsenko 2008 Thm.~3.1, whose
``automorphic correction'' to a BKM superalgebra is not yet
established.

\paragraph{Heal.}
The motivating conjecture, restated in precise form (replacing the
paper's \S1 final sentence):

\begin{conjecture}[Lorgat, motivating, rectified]\label{conj:lorgat-motivating}
Let $X=(S\times E)/(\mathbb Z/N\mathbb Z)$ with
$S$ a $K3$ surface, $\mathbb Z/N\mathbb Z$ acting diagonally via a
symplectic automorphism $g_N\in\mathrm{Aut}_{\mathrm{sympl}}(S)$ of
order~$N$ (Mukai 1988 \textit{Invent.~Math.}~94~Thm.~1.1) and an
$N$-torsion translation on $E$. Let $L\subset\mathrm{Pic}(S)$ be a
lattice polarisation fixed pointwise by $g_N$, and let $h_M=g^s_N$
be an admissible symplectic secondary automorphism of order $M|N$.
Then the Behrend-weighted reduced Donaldson--Thomas $E$-equivariant
zeta function $Z^X_{L,h_M}(q,t,p)$ satisfies
\[
   Z^X_{L,h_M}(q,t,p) \;=\; \frac{C_N}{\Phi^{(N)}(q,t,p)},
   \qquad
   \Phi^{(N)}=\bigl(D^{(N)}_{k(N)}\bigr)^2,
\]
where $D^{(N)}_{k(N)}$ is the unique (up to scaling) Gritsenko--Cl\'ery
diagonal-divisor paramodular cusp form of weight $k(N)$ on
$\Gamma_{t(N)}(N_0)$ for an explicit assignment
$(t(N),N_0(N),k(N),\nu(N))$, and the level-$N$ character of $\Phi^{(N)}$
matches the Borcherds singular theta lift of the $g_N$-twined K3
elliptic genus $\phi^{(g_N)}_{0,1}\in J_{0,1}(\Gamma_0(N))$. Status:

\begin{itemize}[leftmargin=1.5em]
\item $\ClaimStatusTheorem$ at $N=1$: Oberdieck--Pandharipande 2016
  (arXiv:1411.1514) proves $Z^{S\times E}=-1/\Phi_{10}$; the
  reciprocal-square-root side is $\Delta_5$ with
  $\Phi_{10}=\Delta_5^2$; the BKM side is $\mathfrak g_{\Delta_5}$
  per Gritsenko--Nikulin 1998 \textit{J.~Reine Angew.~Math.}~507
  (Thm.~2.1) with denominator identity $\tfrac{1}{64}\Delta_5(2Z)$.
  This is the content of \S\S 2--6 of Lorgat 2020.
\item $\ClaimStatusConjectured$ at $N\in\{2,3,4,6\}$: the identity
  $Z^X_{L,h_M}=C_N/\Phi^{(N)}$ is expected following
  Bryan--Oberdieck--Pandharipande--Yin 2018 and is proved modulo
  the reduced-DT=PT equivalence for the primitive-class sector
  (Oberdieck 2018 \textit{J.~Reine Angew.~Math.}~752, arXiv:1607.01149);
  the BKM automorphic-correction side
  $\mathfrak g_{\Phi^{(N)}}$ is the Gritsenko--Nikulin 1998
  hierarchy at Theorem~2.1 with root multiplicities
  $c_N(\alpha)=\mathrm{ch}^{(g_N)}_{\phi^{(g_N)}_{0,1}}(\alpha)$
  given by $g_N$-twined elliptic-genus Fourier coefficients.
\item $\ClaimStatusOpen$ at $N\in\{5,7,8\}$: the half-integer
  paramodular tower Cl\'ery--Gritsenko 2008 Thm.~3.1 produces
  $D^{(N)}_{k(N)}$ at $k(N)\in\{\tfrac12,1,\tfrac32\}$, but the
  Borcherds-lifted BKM superalgebra at these weights does not
  admit the standard denominator identity and is
  structurally distinct (no integrality of $\kBKM$).
\end{itemize}
\end{conjecture}

%--------------------------------------------------------------------
\subsection*{Cycle 3: Twined elliptic genera -- the $c_N(0)$ constants
and the Mukai--Mathieu arithmetic}
%--------------------------------------------------------------------

\paragraph{Attack.}
The conjecture's third clause: \emph{``root multiplicities
specified by $g_N-h_M$-twisted twined elliptic genera of K3
surfaces.''} Two attacks:

\begin{enumerate}[label=\textbf{O.\arabic*},leftmargin=1.5em]
\item Whose definition of $g_N$-twined elliptic genus? The phrase
  admits two inequivalent conventions:
  \begin{itemize}
    \item Euler-character twining:
      $\chi^{(g_N)}(K3)=\mathrm{tr}_{H^\ast(K3)}(g_N)$, an
      integer; for $g_N$ Mukai-symplectic of order~$N\in\{1,2,3,4,6\}$
      the values are $(24,8,6,4,2)$ (Mukai 1988 Table~1).
    \item Elliptic-genus twining:
      $\phi^{(g_N)}_{0,1}(\tau,z):=\mathrm{tr}_{H^\ast(K3)}
      (g_N(-1)^F y^{J_0}q^{L_0-c/24})$, a weak Jacobi form of
      weight~$0$ and index~$1$ on $\Gamma_0(N)$ (EOT 2011
      \textit{Exp.~Math.}~20; Cheng--Harvey 2012 arXiv:1204.2779
      Thm.~2.3).
  \end{itemize}
  Only the latter is what enters Borcherds' singular theta lift on
  $\Lambda^{3,2}$; the former is its $q^0|_{z=0}$ constant up to a
  factor of~$2$.
\item The Mathieu moonshine connection: Cheng--Duncan--Harvey 2014
  (arXiv:1307.5793, \textit{Commun.~Number Theory
  Phys.}~\textbf{8}) associate to each $M_{24}$ conjugacy class
  $[g]$ a weak Jacobi form $\phi^{[g]}_{0,1}$ via the Mathieu
  moonshine conjecture (proved by Gannon 2012 arXiv:1211.3703 at
  the character level). The Mukai symplectic classes
  $\{1A,2A,3A,4B,6A\}$ of $M_{24}$ correspond precisely to
  $N\in\{1,2,3,4,6\}$ via Mukai's embedding
  $\mathrm{Aut}_{\mathrm{sympl}}(K3)\hookrightarrow M_{23}\subset
  M_{24}$; this is what \emph{makes} $c_N(0)$ well-defined.
  Paper~\S 1 does not cite Mathieu moonshine, reflecting its
  pre-Gannon-era drafting; the content it names is
  Mukai-symplectic twining of elliptic genera.
\end{enumerate}

\paragraph{Ghost theorem.}
For $N\in\{1,2,3,4,6\}$, Mukai 1988 classifies
$\mathrm{Aut}_{\mathrm{sympl}}(K3)$ finite-order elements;
each $g_N$ of order $N\leq 8$ has Euler-characteristic trace
$\chi^{g_N}(K3)\in\{24,8,6,4,2,-,4,0\}$ for
$N\in\{1,2,3,4,6,-,7,8\}$. The twined elliptic genus is
\[
  \phi^{(g_N)}_{0,1}(\tau,z)
  \;=\;
  \mathrm{tr}_{H^{\ast}(K3)}\!\Big(g_N\,(-1)^F y^{J_0}q^{L_0-c/24}\Big),
\]
a weak Jacobi form of weight~$0$, index~$1$, and multiplier system
$\nu^{(g_N)}$ on the Hecke congruence subgroup
$\Gamma_0(N)$. Its $q^0$ Laurent polynomial is
\[
  \phi^{(g_N)}_{0,1}(\tau,z)\big|_{q=0}
  \;=\;
  c_N(0,-1)y^{-1}+c_N(0,0)+c_N(0,1)y,
\]
with $c_N(0,-1)=c_N(0,1)=1$ fixed by the index-$1$ cusp condition,
and $c_N(0,0)\in\{10,4,2,2,2\}$ for $N\in\{1,2,3,4,6\}$
respectively. The scalar
$c_N(0):=c_N(0,0)=\phi^{(g_N)}_{0,1}(\tau,z)\big|_{q=0,z=0}$
equals the $g_N$-twined cohomological Euler character of $K3$
normalised by the $y^0$-coefficient; its values match those of
Cheng--Harvey 2012 Table~1 at $N\in\{1,2,3,4,6\}$.

The Borcherds singular theta lift on $\Lambda^{3,2}$ with input
$\phi^{(g_N)}_{0,1}$ yields a level-$N$ Siegel paramodular cusp form
$\Phi^{(N)}=\mathrm{Lift}(\phi^{(g_N)}_{0,1})$ of weight
$\mathrm{wt}(\Phi^{(N)})=c_N(0)/2\in\{5,2,1,1,1\}$
(Borcherds 1998 \textit{Invent.~Math.}~132~Thm.~13.3). Matching the
Gritsenko--Nikulin 1998 $\Delta^{(N)}_{k(N)}$-tower identifies
\[
  (k(N))_{N\in\{1,2,3,4,6\}}=(5,2,1,1,1)=(c_N(0)/2)_{N\in\{1,2,3,4,6\}},
\]
so
\[
  \boxed{\kBKM(\Phi^{(N)}) \;=\; c_N(0)/2, \quad
    N\in\{1,2,3,4,6\}.}
\]
At $N=1$ this recovers $\kBKM(\Delta_5)=10/2=5$ as the
Gelfand-check of Lorgat 2020's central numerical assertion.

\paragraph{Heal.}
The twining-side of the rectified conjecture
(Conjecture~\ref{conj:lorgat-motivating}) is inscribed as:

\begin{proposition}[Twined Borcherds input with $c_N(0)$-table]
\label{prop:twined-cN}
For $N\in\{1,2,3,4,6\}$ and $g_N\in\mathrm{Aut}_{\mathrm{sympl}}(K3)$
Mukai-symplectic of order $N$:
\begin{center}
\begin{tabular}{cccccc}
\hline
$N$ & $[g]\in M_{24}$ & $\chi^{g_N}(K3)$ & $c_N(0)$ &
$k(N)=c_N(0)/2$ & $D^{(N)}_{k(N)}$ \\
\hline
$1$ & $1A$ & $24$ & $10$ & $5$ & $\Delta_5$ \\
$2$ & $2A$ & $8$ & $4$ & $2$ & $D^{(2)}_2$ \\
$3$ & $3A$ & $6$ & $2$ & $1$ & $D^{(3)}_1$ \\
$4$ & $4B$ & $4$ & $2$ & $1$ & $D^{(4)}_1$ \\
$6$ & $6A$ & $2$ & $2$ & $1$ & $D^{(6)}_1$ \\
\hline
\end{tabular}
\end{center}
The table is an identity of Fourier-coefficient arithmetic between
(a) the Mukai-symplectic classification of $\mathrm{Aut}_{\mathrm{sympl}}(K3)$,
(b) the $M_{24}$ twined elliptic genus, and (c) the
Gritsenko--Cl\'ery paramodular tower. Primary: Mukai 1988;
Cheng--Harvey 2012; Gritsenko--Cl\'ery 2008; Borcherds 1998.
Proof: direct computation of $\phi^{(g_N)}_{0,1}|_{q=0,z=0}$ and
Borcherds Thm.~13.3.
\end{proposition}

The $\kBKM$-subscript discipline forbids writing ``$\kBKM=c(0)/2$''
without the $N$-index on both sides. The rectified motivating
conjecture's twining clause reads:
\emph{``root multiplicities $\mult_{\Phi^{(N)}}(\alpha)=c_N(nm,\ell)$
where $\alpha=(n,\ell,m)\in\Lambda^{2,1}_{II}\cap\mathbb R_{>0}
\mathcal P_{II}$ and $c_N(nm,\ell)$ is the $(nm,\ell)$-Fourier
coefficient of the $g_N$-twined K3 elliptic genus
$\phi^{(g_N)}_{0,1}$.''}

%--------------------------------------------------------------------
\subsection*{Cycle 4: The Delta\_5 weight is 5, not 10 -- Borcherds
1998 Thm.~13.3 and the Siegel/paramodular discipline}
%--------------------------------------------------------------------

\paragraph{Attack.}
The paper's~\S 2 opens: \emph{``$\Delta_{10}$ is the square of a cusp
form $\Delta_5$ of weight~$5$.''} Attacks:

\begin{enumerate}[label=\textbf{O.\arabic*},leftmargin=1.5em]
\item Which convention for ``weight''? A Siegel modular form of
  weight $k$ on $\Sp_4(\mathbb Z)$ transforms as $F(MZ)=(\det
  CZ+D)^k\nu(M)F(Z)$; a paramodular form of weight $k$ on
  $\Gamma_t$ at level $t$ transforms by the same formula on
  $\Gamma_t$. The numerical weight on
  $\Gamma_1=\Sp_4(\mathbb Z)$ is a single number; but
  $\Phi_{10}\in M_{10}(\Sp_4(\mathbb Z))$ and
  $\Delta_5\in M_5(\Sp_4(\mathbb Z),\nu_{\Delta_5})$ (with
  non-trivial character of order~$2$) are \emph{both} written on the
  same group. The ambiguity is not between Siegel and paramodular;
  it is between $\Phi_{10}$ and $\Delta_5$, and it costs a factor
  of~$2$.
\item What is the Borcherds-lift weight? Borcherds 1998
  \textit{Invent.~Math.}~132 Theorem~13.3 (``Automorphic forms on
  $O_{s+2,2}(\mathbb R)$ and infinite products'', primary citation)
  states: if
  $\phi(\tau)=\sum_{\lambda}q^{-n_\lambda}c_\lambda(-n_\lambda)e_\lambda
  +\cdots$ is a weakly holomorphic vector-valued modular form of
  weight $1-(n^++n^-)/2$ on a lattice $\Lambda$ of signature
  $(n^+,n^-)$ (the bookkeeping lift to vector-valued Jacobi forms
  proceeds via Eichler--Zagier 1985 Thm.~5.1), then its singular
  theta lift
  $\Phi(Z) :=\lim_{\epsilon\to 0^+}\mathrm{reg}\int_{\Gamma_1
    \backslash\mathfrak H_1}\Theta_\Lambda(\tau,Z)\phi(\tau)
  y^{-\epsilon}\,\frac{dxdy}{y^2}$
  is an automorphic form on $\mathrm O(\Lambda)_+$ of weight
  $\mathrm{wt}(\Phi)=c_0(0)/2$, with $c_0(0)$ the
  $q^0$ coefficient of the trivial-class component.
\item At $\Lambda=\Lambda^{3,2}$ (signature $(3,2)$) and
  $\phi=\phi_{0,1}$ (weight~$0$, index~$1$, lifted to
  vector-valued form of weight $1-5/2=-3/2$ on $\Lambda^{3,2}$ via
  the Eichler--Zagier $\theta$-decomposition), the $q^0$ constant
  of the trivial-class component is $c_{\phi_{0,1}}(0,0)$, the
  $y^0$-coefficient of $\phi_{0,1}(\tau,z)|_{q=0}=y^{-1}+10+y$. So
  $c_{\phi_{0,1}}(0,0)=10$, and Borcherds' formula gives
  $\mathrm{wt}(\Phi)=10/2=5$, not~$10$.
\end{enumerate}

\paragraph{Ghost theorem.}
The Borcherds singular theta lift on the quadratic lattice
$\Lambda=\Lambda^{3,2}\simeq\Lambda^{1,1}\oplus\Lambda^{1,1}\oplus[2]$
with input the weight-$0$, index-$1$ weak Jacobi form
$\phi_{0,1}(\tau,z)$ has weight
\[
   \mathrm{wt}\bigl(\mathrm{Lift}(\phi_{0,1})\bigr)
   \;=\;\frac{c_{\phi_{0,1}}(0,0)}{2}\;=\;\frac{10}{2}\;=\;5
\]
by Borcherds 1998 \textit{Invent.~Math.}~\textbf{132}~Thm.~13.3. The
explicit output
\[
   \mathrm{Lift}(\phi_{0,1})(Z) \;=\; \frac{1}{64}\,\Delta_5(2Z)
\]
is the content of Lorgat 2020 Theorem~3
(\texttt{wave13\_source\_paper.txt:689--693}): Gritsenko--Nikulin
1995 Part~II (arXiv:alg-geom/9504006) Theorem~2.1. The factor
$\tfrac{1}{64}$ is the $f(1,1,1)$-normalisation (cycle~5).

The Igusa cusp form $\Phi_{10}\in M_{10}(\Sp_4(\mathbb Z))$ of weight
$10$ is then $\Phi_{10}=\Delta_5^2$ with weight~$10$, character
$\nu_{\Delta_5}^2=\mathbf 1$ trivial; both $\Delta_5$ and $\Phi_{10}$
live on $\Sp_4(\mathbb Z)=\Gamma_1$ and both are well-defined Siegel
modular forms. The Gelfand-level numerical match is
\[
  \kBKM(\Delta_5) \;=\; \mathrm{wt}(\Delta_5)\;=\;5
  \;=\;c_{\phi_{0,1}}(0,0)/2,
\]
with $c_{\phi_{0,1}}(0,0)=10$ the $q^0|_{z=0}$ constant of the weak
Jacobi form input.

\paragraph{Heal.}
The Delta\_5 preamble (rectified):

\begin{definition}[Igusa cusp form and its square root]
\label{def:delta-5-preamble}
The ring of Siegel modular forms of genus~$2$ with respect to the
full Siegel modular group $\Sp_4(\mathbb Z)$ is
$\mathrm{SM}(\Sp_4(\mathbb Z))=\mathbb C[E_4,E_6,\Delta_{10},\Delta_{12}]$
(Igusa 1964 \textit{Am.~J.~Math.}~\textbf{86}), with $E_4,E_6$ the
genus-$2$ Eisenstein series and $\Delta_{10}\in M_{10}(\Sp_4(\mathbb Z)),
\Delta_{12}\in M_{12}(\Sp_4(\mathbb Z))$ the two cusp-form
generators of weights~$10,12$.

The Igusa cusp form $\Phi_{10}:=\Delta_{10}$ is a square: there
exists a Siegel cusp form
\[
   \Delta_5 \in M_5\bigl(\Sp_4(\mathbb Z),\,\nu_{\Delta_5}\bigr)
\]
of weight~$5$ and non-trivial multiplier system $\nu_{\Delta_5}$ of
order~$2$ such that $\Phi_{10}=\Delta_5^2$ (Igusa 1964; the
character was determined explicitly by Maass 1964
\textit{Nachr.~Akad.~Wiss.~G\"ottingen Math.-Phys.~Kl.}~11).
Explicitly,
\[
   \Delta_5 \;=\; \prod_{\substack{(a,b)\in(\mathbb Z/2\mathbb Z)^2\\
     {}^t\! a b\equiv 0\pmod 2}}
     \vartheta_{a,b},
\]
the product of the ten even theta constants on the genus-$2$
theta-characteristics $(a,b)\in(\mathbb Z/2\mathbb Z)^4$ with
$a\cdot b\equiv 0\pmod 2$. The multiplier system $\nu_{\Delta_5}$
has explicit formulae on the generators of $\Sp_4(\mathbb Z)$
given by Maass 1964; in the shearing decomposition
$M=\mathrm{diag}({}^t\! A^{-1},A)\cdot\bigl(\begin{smallmatrix}
I_2 & B\\0 & I_2\end{smallmatrix}\bigr)\cdot
\bigl(\begin{smallmatrix}0 & I_2\\-I_2 & 0\end{smallmatrix}\bigr)$,
\begin{align*}
  \nu_{\Delta_5}\!\Big(\begin{smallmatrix}0 & I_2\\-I_2 & 0\end{smallmatrix}\Big)
  &=1,
  \\
  \nu_{\Delta_5}\!\Big(\begin{smallmatrix}I_2 & B\\0 & I_2\end{smallmatrix}\Big)
  &=(-1)^{b_1+b_2+b_3},
  \\
  \nu_{\Delta_5}\!\Big(\begin{smallmatrix}{}^t\! A^{-1} & 0\\0 & A\end{smallmatrix}\Big)
  &=(-1)^{(1+a_1+a_4)(1+a_2+a_3)+a_1a_4}.
\end{align*}
\end{definition}

\begin{theorem}[Borcherds lift of $\phi_{0,1}$ to $\Delta_5$]
\label{thm:borcherds-phi01-delta5}
Let $\phi_{0,1}(\tau,z)=\sum_{n\geq 0,\ell\in\mathbb Z}
f(n,\ell)q^n y^\ell$ be the unique-up-to-scaling weak Jacobi form
of weight~$0$ and index~$1$ (Eichler--Zagier 1985
\textit{The Theory of Jacobi Forms}, Prop.~2.1; normalisation
$\phi_{0,1}|_{q=0}=y^{-1}+10+y$, i.e.\ $f(0,\pm1)=1$ and $f(0,0)=10$).
Then the Borcherds singular theta lift on
$\Lambda=\Lambda^{3,2}\simeq\Lambda^{1,1}\oplus\Lambda^{1,1}\oplus[2]$
with input $\phi_{0,1}$ (decomposed to vector-valued form on
$\Lambda^{3,2}$ via Eichler--Zagier $\theta$-decomposition) produces
the Siegel modular form
\[
  \mathrm{Lift}(\phi_{0,1})(Z) \;=\; \tfrac{1}{64}\,\Delta_5(2Z)
\]
of weight $\mathrm{wt}(\Delta_5)=f(0,0)/2=5$, matching Borcherds
1998 \textit{Invent.~Math.}~132 Thm.~13.3. The GKM superalgebra
$\mathfrak g_{\Delta_5}$ constructed in \S 5 of Lorgat 2020 has
Weyl--Kac--Borcherds denominator
\[
   \Phi_{\mathfrak g_{\Delta_5}}(z) \;=\; \tfrac{1}{64}\,\Delta_5(2z)
\]
(Gritsenko--Nikulin 1995 Part~II Theorem~2.1; Lorgat 2020 Theorem~3).
Status: $\ClaimStatusTheorem$, with the decomposition
$\Delta_5\mapsto\phi_{0,1}$ the content of the paper's \S 2 Lemma and
\S 6 Theorem~4.
\end{theorem}

%--------------------------------------------------------------------
\subsection*{Cycle 5: The $f(1,1,1)=64$ normalisation and the infinite
product -- what the 64 is, and what ``product over positive roots'' means}
%--------------------------------------------------------------------

\paragraph{Attack.}
The paper's~\S 2 states: \emph{``$f(1,1,1)=64$ as well as
$64\mid f(n,\ell,m)$ for all $n,\ell,m$''} and uses this without
explaining why the $64$ factor arises or how it propagates. Three
attacks:

\begin{enumerate}[label=\textbf{O.\arabic*},leftmargin=1.5em]
\item The factor $64=2^6$ has an immediate arithmetic origin: $2^6$
  is the size of the kernel of the twofold cover
  $\mathrm{Sp}_4(\mathbb Z)\to\mathrm{Sp}_4(\mathbb Z)/\{\pm I\}$
  tensored with the genus-$2$ theta-characteristic structure
  (ten even theta constants at first nontrivial Fourier index each
  contribute a factor of $2$: but this heuristic alone gives
  $2^{10}$, not $2^6$). So the $64$ is not ``obvious from the theta
  constants''; it demands a computation.
\item Why does $64\mid f(n,\ell,m)$ for all $n,\ell,m$? The Lorgat
  2020 derivation proceeds via the Fourier--Jacobi
  expansion
  \[
    \Delta_5(Z)=\sum_{m>0,m\equiv 1\pmod 2}\phi_{5,m}(z_1,z_2)
      \,\exp(\pi imz_3),
  \]
  where $\phi_{5,m}$ is a Jacobi cusp form of weight~$5$ and
  index~$\tfrac m2$, and the first Fourier--Jacobi coefficient
  factors as $\phi_{5,\tfrac12}=\eta(z_1)^9\vartheta_{11}(z_1,z_2)$,
  then squares to a Jacobi cusp form of weight~$10$ index~$1$; the
  $64$ is the product of normalisation factors $\eta^9,
  \vartheta_{11}$ at the first nontrivial Fourier mode, as Lorgat
  2020 \S 2 Lemma computes.
\item The infinite-product expansion of Lorgat 2020 Theorem~4
  \[
    \tfrac{1}{64}\Delta_5(Z)
    = \exp(\pi i(z_1+z_2+z_3))\prod_{(n,\ell,m)>0}
      \bigl(1-\exp(2\pi i(nz_1+\ell z_2+mz_3))\bigr)^{f(nm,\ell)}
  \]
  depends on: (i) the cone-ordering ``$(n,\ell,m)>0$''
  (defined in Lorgat 2020 Remark~1: $n,m\geq 0$, $\ell$ arbitrary
  if $n>0$ or $m>0$, $\ell<0$ if $n=m=0$); (ii) the identification
  of multiplicities $\mult(\alpha=(n,\ell,m))=f(nm,\ell)$ with
  $\phi_{0,1}$-Fourier coefficients; and (iii) the Borcherds
  singular-theta-lift cone-ordering, which differs from the
  Weyl-chamber ordering by the sequence of cone embeddings
  $\Delta(\Lambda^{2,1}_{II})^*_+\subset\mathcal C(\Lambda^{2,1})_+
  =\mathcal C(\Lambda^{2,1})^*_+\subset\Delta(\Lambda^{2,1}_{II})_+$
  (Lorgat 2020 \S 4).
\end{enumerate}

\paragraph{Ghost theorem.}
The infinite-product expansion is the Borcherds--Gritsenko--Nikulin
automorphic-correction formula. Expressed cleanly:

\begin{theorem}[Infinite-product expansion of $\Delta_5$; Lorgat 2020
Theorem~4]\label{thm:delta5-infinite-product}
Let $\phi_{0,1}(\tau,z)=\sum_{n,\ell}f(n,\ell)q^n y^\ell$ with
$f(n,\ell)$ dependent only on $4n-\ell^2$ (Eichler--Zagier 1985
Thm.~2.2, parity of weight). Then
\[
  \tfrac{1}{64}\Delta_5(Z)
  \;=\;
  \exp\!\bigl(\pi i(z_1+z_2+z_3)\bigr)\!
  \prod_{(n,\ell,m)\in\mathbb Z^3,\ (n,\ell,m)>0}\!\!
  \bigl(1-\exp(2\pi i(nz_1+\ell z_2+mz_3))\bigr)^{f(nm,\ell)}
\]
where the positive-root cone $(n,\ell,m)>0$ is
\[
   (n,\ell,m)>0
   \;\Longleftrightarrow\;
   \bigl(n\geq 0,\ m\geq 0,\ (n,m)\neq(0,0);\ \ell\in\mathbb Z\bigr)
   \;\vee\;
   \bigl(n=m=0,\ \ell<0\bigr).
\]
The multiplicity function $\mult(\alpha)=f(nm,\ell)$ is the
Fourier coefficient of $\phi_{0,1}$, equivalently the
super-dimension of the root space
$(\mathfrak g_{\Delta_5})_{\alpha=(n,\ell,m)}$ in the
$\Lambda^{2,1}_{II}$-grading. Proof: Fourier--Jacobi expansion of
$\Delta_5$, Borcherds singular theta lift on $\Lambda^{3,2}$, and
the Weyl--Kac--Borcherds denominator identity for
$\mathfrak g_{\Delta_5}$ applied to the trivial one-dimensional
representation.
\end{theorem}

\paragraph{Heal.}
The preamble's $f(1,1,1)=64$ line is replaced by:

\begin{lemma}[Normalisation of the first Fourier coefficient]
\label{lem:f111-64}
In the Fourier expansion
$\Delta_5(Z)=\sum_{n,\ell,m=1\pmod 2;\ 4nm-\ell^2>0;\ n,m>0}
f(n,\ell,m)\,\exp(\pi i(nz_1+\ell z_2+mz_3))$,
the leading coefficient is $f(1,1,1)=64$, and $64\mid f(n,\ell,m)$
for all admissible $(n,\ell,m)$. Equivalently, the first
Fourier--Jacobi coefficient factors as
$\phi_{5,\tfrac12}(z_1,z_2)=\eta(z_1)^9\vartheta_{11}(z_1,z_2)$, and
\[
  \tfrac{1}{64}\phi_{5,1}(z_1,z_2)
  \;=\;
  \psi_{5,\tfrac12}(z_1,z_2)^2
  \;=\;
  -q^{1/2}y^{-1/2}\prod_{n\geq 1}(1-q^{n-1}y)(1-q^n y^{-1})(1-q^n)^{10},
\]
the unique (up to scaling) Jacobi cusp form of weight~$10$ and
index~$1$, identified with the first Fourier--Jacobi coefficient of
$\Delta_{10}=\Delta_5^2$ (Lorgat 2020 \S 2 Lemma; Jacobi triple
product applied to the $y^{1/2}$ coefficient).
\end{lemma}

%--------------------------------------------------------------------
\subsection*{Convergence: the motivating conjecture in rectified form,
Chriss--Ginzburg economy}
%--------------------------------------------------------------------

The five cycles close. The rectified \S 1--\S 2 reads, in the voice
of Gelfand (direct construction, no hedging, no narration):

\begin{quote}
\emph{Fix genus~$g=2$ and consider the paramodular groups
$\Gamma_t(N)<\Sp_4(\mathbb Q)$ of Hecke type at level
$(t,N)\in\mathbb Z_+^2$. Gritsenko--Cl\'ery 2008 classify the
diagonal-divisor Siegel paramodular cusp forms: the one-dimensional
spaces $\mathbb C\cdot D^{(N)}_{k(N)}$ for the eight triples
$(t,N,k,\nu)$ listed in Proposition~\ref{prop:twined-cN}; the
archetype is $D^{(1)}_5=\Delta_5\in M_5(\Sp_4(\mathbb Z),\nu_{\Delta_5})$,
the unique square root of the Igusa cusp form
$\Phi_{10}=\Delta_5^2$.}

\emph{Each $D^{(N)}_{k(N)}$ is the Borcherds singular theta lift
(\textit{Invent.~Math.}~132 Thm.~13.3) of the $g_N$-twined K3
elliptic genus $\phi^{(g_N)}_{0,1}\in J_{0,1}(\Gamma_0(N))$, with
weight $k(N)=c_N(0)/2=\mathrm{tr}_{H^*(K3)}(g_N)|_{q=0,z=0}/2$ and
multiplier character determined by the $\Gamma_0(N)$-level of the
twisting automorphism. The BKM superalgebra automorphic correction
$\mathfrak g_{\Phi^{(N)}}$ (Gritsenko--Nikulin 1998
\textit{J.~Reine Angew.~Math.}~507 Thm.~2.1) has denominator
identity $\Phi^{(N)}(Z)=(D^{(N)}_{k(N)}(Z))^2$ and root
multiplicities $\mult(\alpha)=c_N(nm,\ell)$ given by the
$\phi^{(g_N)}_{0,1}$-Fourier coefficients.}

\emph{On the arithmetic geometry side, the Behrend-weighted reduced
$E$-equivariant Donaldson--Thomas zeta $Z^X_{L,h_M}$ of the
CHL-orbifold $X=(K3\times E)/(\mathbb Z/N\mathbb Z)$ satisfies the
Igusa-cusp-form reciprocal identity
$Z^X_{L,h_M}=C_N/\Phi^{(N)}=C_N/(D^{(N)}_{k(N)})^2$. Status:
$\ClaimStatusTheorem$ at $N=1$ via Oberdieck--Pandharipande 2016
\textit{Invent.~Math.}~213 and Oberdieck--Pixton 2018
\textit{Duke Math.~J.}~168; $\ClaimStatusConjectured$ at
$N\in\{2,3,4,6\}$ per Bryan--Oberdieck--Pandharipande--Yin 2018.
The three correspondences ---
(Gritsenko--Cl\'ery classification) $\leftrightarrow$
(twined K3 elliptic genera of Mukai-symplectic automorphisms)
$\leftrightarrow$ (reciprocal-square-roots of reduced
CHL-orbifold DT zeta functions) --- are the content of the
motivating conjecture.}
\end{quote}

Six routes of construction in Vol~III (CoHA, stable envelope,
Koszul, chiral Yangian via $\Phi$, BKM screening, Borcherds lift)
are six different constructions, not six applications of
$\Phi_3^{\mathrm{FA}}$ to $D^b(\mathrm{Coh}(K3\times E))$
(cache~F8, AP-CY144). The specialisation
$\mathrm{Sp}_{K3,E}(\Phi_3^{\mathrm{FA}}(K3\times E))$ of
Wave~12~U1 is the canonical chiral shadow; $\mathfrak g_{\Delta_5}$
sits on the classical-limit BPS Lie subalgebra of this shadow.

%--------------------------------------------------------------------
\subsection*{Two-stage factorisation perspective (cross-reference to
Wave~12 U1 and cache F8)}
%--------------------------------------------------------------------

In the two-stage factorisation of Wave~12~U1
(\texttt{notes/wave12\_u1\_two\_stage\_functor.tex}),
\[
  \Phi_d \;=\; \mathrm{Sp}_{\Sigma_{d-1},C}\;\circ\;\Phih_d,
\]
Lorgat 2020 is the explicit computation of the specialisation
$\mathrm{Sp}_{K3,E}$ applied to the native $E_3$-holomorphic
factorisation algebra $\Phih_3(K3\times E)$:

\begin{itemize}[leftmargin=1.5em]
\item Stage~1 canonical: $\Phih_3(K3\times E)$ is an $E_3$-holomorphic
  factorisation algebra on the compact CY$_3$
  $X=K3\times E$, the Costello--Li 2020 hCS BV construction on $X$
  (arXiv:1505.06703 \S 3).
\item Stage~2 choice-dependent: the fibre cycle $\Sigma_2=K3$ with
  reference curve $C=E$ gives the $E_1$-chiral shadow
  $\mathrm{Sp}_{K3,E}(\Phih_3(K3\times E))$ on the reference
  elliptic curve $E$. This is a chiral bialgebra whose
  classical-limit BPS Lie subalgebra is conjecturally
  $\mathfrak g_{\Delta_5}$.
\item Lorgat 2020 constructs $\mathfrak g_{\Delta_5}$ directly from
  the Borcherds lift of $\phi_{0,1}$ on $\Lambda^{3,2}$, without
  passing through the factorisation-algebra formalism; the result
  matches the classical-limit BPS Lie subalgebra of
  $\mathrm{Sp}_{K3,E}(\Phih_3)$ at the denominator-identity level
  (Gritsenko--Nikulin 1998 matches Borcherds 1998 via the common
  $\Lambda^{3,2}$-lift).
\item Sibling specialisations $\mathrm{Sp}_{\Sigma'_2,C}$ for other
  cycles $\Sigma'_2\in\mathrm{Cyc}_2(K3\times E)$ give different
  BKMs: the $\mathrm{Sp}_{\Sigma^{\mathrm{Fake}}_2,E}$ specialisation
  gives a Fake-Monster sibling; the $\mathrm{Sp}_{H_N,E}$
  specialisation for a paramodular Humbert surface $H_N$ gives the
  sibling $\Phi^{(N)}$-BKM at level $N$. Lorgat 2020's
  $N\in\{2,3,4,6\}$ cases are the $\mathrm{Sp}_{H_N,E}$-family.
\end{itemize}

The motivating conjecture of Lorgat 2020 thus translates under the
two-stage factorisation to: \emph{the eight Gritsenko--Cl\'ery
diagonal-divisor forms are the weight-$c_N(0)/2$ denominators of eight
distinct $\mathrm{Sp}_{\Sigma_2,C}$-specialisations of
$\Phih_3$ applied to CHL-orbifolds of $K3\times E$
(and their $N=7,8$ half-integer analogues on
$\Phih_3$ of $N$-twisted CY$_3$s).}

%--------------------------------------------------------------------
\subsection*{Status summary, survived / falsified / new derivations}
%--------------------------------------------------------------------

\begin{description}
\item[\ClaimStatusTheorem ~Survived in strong form]
Lorgat 2020 Theorem~2 ($Z^{S\times E}=C/(\Delta_5)^2$ at $N=1$):
is a theorem, not a conjecture, via
Oberdieck--Pandharipande 2016. Lorgat 2020 Theorem~3
($\tfrac{1}{64}\Delta_5(2Z)=\Phi(z)$ as
Weyl--Kac--Borcherds denominator of $\mathfrak g_{\Delta_5}$):
theorem via Gritsenko--Nikulin 1998 Thm.~2.1.
Lorgat 2020 Theorem~4 (infinite product for
$\tfrac{1}{64}\Delta_5$): theorem via the same source.
$\kBKM(\Delta_5)=5=c_{\phi_{0,1}}(0,0)/2$: theorem via
Borcherds 1998 Thm.~13.3 (not Thm.~10.1).

\item[\ClaimStatusConjectured ~Survived as conjecture]
Lorgat 2020 Conjecture~1 (generalisation to
$N\in\{2,3,4,6\}$ CHL-orbifolds): conjectural, partially proved
for the reduced-primitive-class sector via Oberdieck 2018
\textit{J.~Reine Angew.~Math.}~752 and
Bryan--Oberdieck--Pandharipande--Yin 2018. Lorgat 2020
Conjecture~1 at $N\in\{5,7,8\}$ (half-integer paramodular tower):
open; requires Cl\'ery--Gritsenko 2008 half-integer Borcherds
lift, whose BKM automorphic-correction is not yet constructed.

\item[\ClaimStatusDerivation ~New derivations inscribed]
\begin{enumerate}[leftmargin=1.5em]
  \item The eight Gritsenko--Cl\'ery forms enumerated with explicit
    $(t,N,k,\nu)$-signature (cycle~1 heal table).
  \item The $c_N(0)$ table at $N\in\{1,2,3,4,6\}$:
    $c_N(0)=(10,4,2,2,2)$, matched against Mukai-symplectic
    orders (Proposition~\ref{prop:twined-cN}).
  \item The weight-5 derivation via Borcherds 1998 Thm.~13.3:
    $\mathrm{wt}(\Delta_5)=f(0,0)/2=5$, where $f(0,0)=10$ is the
    $y^0$-coefficient of $\phi_{0,1}|_{q=0}=y^{-1}+10+y$
    (Theorem~\ref{thm:borcherds-phi01-delta5}).
  \item The $f(1,1,1)=64$ normalisation as a product of
    $\eta^9\vartheta_{11}$ factors at the first nontrivial
    Fourier mode (Lemma~\ref{lem:f111-64}).
\end{enumerate}

\item[Falsified]
The unsubscripted-$\kBKM$ / unsubscripted-$c(0)$ phrasings (``$\kBKM=5=c(0)/2$''
without $N$-subscript) of earlier Vol~III drafts: subscript
discipline ($\kBKM(\Phi^{(N)})=c_N(0)/2$ with $N$ explicit on both
sides). Not a falsification of content; a falsification of
convention. The Kazhdan-voice Wave~12~A2 assault already
enforces this at the cache-AP113 level.

The claim that \emph{``all eight [Gritsenko--Cl\'ery] specimens arise
as Borcherds lifts''} is not fully established at the half-integer
tower ($N\in\{4,7,8\}$ weight-$\tfrac12,\tfrac32$ specimens):
Cl\'ery--Gritsenko 2008 Thm.~3.1 constructs these half-integer
paramodular forms directly (not via Borcherds lift of a Jacobi
form), and the lift-to-BKM correspondence fails integrality. Three
of the eight specimens cross an integrality boundary; the
motivating conjecture as stated must be qualified to
$N\in\{1,2,3,4,6\}$ for the BKM-lifting clause to hold.
\end{description}

%--------------------------------------------------------------------
\subsection*{Frontier: open questions for Wave~14}
%--------------------------------------------------------------------

\begin{enumerate}[leftmargin=1.5em]
\item \textbf{BKM-side of half-integer weight forms.} Construct the
  BKM automorphic correction for the three half-integer specimens
  $(D^{(4)}_{1/2},D^{(1,4)}_{3/2},D^{(2,2)}_1)$, or prove no such
  correction exists and the three specimens fall outside the
  Borcherds-lift programme. Primary path: Cl\'ery--Gritsenko 2008
  arXiv:0812.3962 \S 4 (half-integer Jacobi input);
  Borcherds 1998 Thm.~13.3 (integer-weight hypothesis) extension
  via Rayen--Shelton 2018 arXiv:1810.03893 (meromorphic Borcherds
  lifts).
\item \textbf{Oberdieck--Pandharipande--Yin completion for
  $N\in\{5,7,8\}$.} Extend the CHL Igusa conjecture; the obstruction
  is the failure of $\chi^{g_N}(K3)\in\{0,4\}$ at these orders to
  produce an integer $c_N(0)/2$ weight.
\item \textbf{Explicit $\mathrm{Sp}_{\Sigma_2,C}$ at the eight
  specimens.} Identify for each of the eight Gritsenko--Cl\'ery
  specimens the specialisation cycle $\Sigma_2\in\mathrm{Cyc}_2$
  that produces it as $\mathrm{Sp}_{\Sigma_2,C}\circ\Phih_3$. At
  $N=1$ this is the $K3$-fibre; at $N=2,3,4,6$ the paramodular
  Humbert surface $H_N$; at half-integer weight the cycle is
  expected to be a lattice-polarised $K3$-fibre with
  $g_N$-invariant polarisation.
\item \textbf{Chiral-side upgrade of $\mathfrak g_{\Delta_5}$.} The
  classical-limit BPS Lie subalgebra of
  $\mathrm{Sp}_{K3,E}(\Phih_3(K3\times E))$ is conjecturally
  $\mathfrak g_{\Delta_5}$; upgrade this to a chiral-bialgebra
  statement at the level of $\mathbf H_{\Delta_5}$ (Wave 15--16
  DRINFELD/GELFAND cache appends), not merely at the Lie-algebra
  bracket level.
\item \textbf{Mathieu moonshine at the BKM level.} Cheng--Duncan--Harvey
  2014 Umbral Moonshine conjecturally extends the Mukai-symplectic
  $c_N(0)$-table to the full $M_{24}$-equivariant $c_{[g]}$-table
  for $[g]$ any $M_{24}$ conjugacy class; extend the rectified
  motivating conjecture to the Umbral regime by identifying the
  $M_{24}$-twined $\Phi^{[g]}$ as Borcherds lifts of
  $\phi^{[g]}_{0,1}$ and their BKM automorphic corrections as
  $\mathfrak g_{\Phi^{[g]}}$.
\end{enumerate}

%--------------------------------------------------------------------
\subsection*{Primary sources audit}
%--------------------------------------------------------------------

\begin{itemize}[leftmargin=1.5em]
\item Igusa 1964 \textit{Am.~J.~Math.}~\textbf{86}: ring of Siegel
  modular forms $\mathrm{SM}(\Sp_4(\mathbb Z))=\mathbb C[E_4,E_6,
  \Delta_{10},\Delta_{12}]$.
\item Maass 1964 \textit{Nachr.~Akad.~Wiss.~G\"ottingen Math.-Phys.
  Kl.}~11: the $\nu_{\Delta_5}$ character of order~$2$.
\item Eichler--Zagier 1985 \textit{The Theory of Jacobi Forms},
  Prog.~Math.~\textbf{55}, Thm.~5.1: decomposition of Jacobi forms into
  vector-valued modular forms, and the uniqueness of $\phi_{0,1}$
  up to scaling.
\item Borcherds 1995 \textit{Invent.~Math.}~\textbf{120} \S 6:
  automorphic forms with singularities on Grassmannians,
  denominator formulae.
\item Borcherds 1998 \textit{Invent.~Math.}~\textbf{132} Thm.~13.3:
  weight of the singular theta lift, $\mathrm{wt}(\Phi)=c_0(0)/2$.
\item Gritsenko--Nikulin 1995 Part~II arXiv:alg-geom/9504006 Thm.~2.1:
  Siegel automorphic correction of Lorentzian Kac--Moody Lie
  algebras; $\tfrac{1}{64}\Delta_5(2Z)=\Phi_{\mathfrak g_{\Delta_5}}$.
\item Gritsenko--Nikulin 1998 \textit{J.~Reine Angew.~Math.}~\textbf{507}:
  BKM hierarchy at $N\in\{1,2,3,4,6\}$, weights $(5,2,1,1,1)$.
\item Gritsenko--Cl\'ery 2008 arXiv:0812.3962: classification of
  $8$ diagonal-divisor paramodular cusp forms at $(t,N,k,\nu)$.
\item Cl\'ery--Gritsenko 2008 Thm.~3.1: half-integer paramodular
  tower at $N=4$, weights $1/2, 3/2, 1$.
\item Mukai 1988 \textit{Invent.~Math.}~\textbf{94} Thm.~1.1:
  $\mathrm{Aut}_{\mathrm{sympl}}(K3)\hookrightarrow M_{23}\subset M_{24}$,
  order at most $8$, Euler-character table at orders
  $\{1,2,3,4,5,6,7,8\}$.
\item EOT 2011 \textit{Exp.~Math.}~\textbf{20}: K3 elliptic genus
  and the $M_{24}$ Mathieu moonshine observation.
\item Cheng--Harvey 2012 arXiv:1204.2779 Thm.~2.3: twined K3 elliptic
  genera $\phi^{[g]}_{0,1}\in J_{0,1}(\Gamma_0(N_g))$ at the $M_{24}$
  conjugacy classes; Table~1: $c_{[g]}(0)$ values.
\item Oberdieck--Pandharipande 2016 arXiv:1411.1514
  \textit{Invent.~Math.}~\textbf{213}: $Z^{S\times E}=-1/\Phi_{10}$
  at $N=1$.
\item Oberdieck 2018 \textit{J.~Reine Angew.~Math.}~\textbf{752}
  (arXiv:1607.01149): reduced DT = PT for CHL-orbifolds of $K3\times E$
  at $N\in\{2,3,4,6\}$, primitive-class sector.
\item Oberdieck--Pixton 2018 \textit{Duke Math.~J.}~\textbf{168}
  (arXiv:1706.10100): Igusa cusp form conjecture at full
  cohomology, holomorphic anomaly equation.
\item Bryan--Oberdieck--Pandharipande--Yin 2018: CHL-orbifold
  extension of the Igusa cusp form conjecture.
\item Cheng--Duncan--Harvey 2014 arXiv:1307.5793
  \textit{Commun.~Number Theory Phys.}~\textbf{8}: Umbral Moonshine.
\item Gannon 2012 arXiv:1211.3703: Mathieu moonshine proved at the
  character level.
\item Costello--Li 2020 arXiv:1505.06703 \S 3: hCS BV construction on
  CY$_3$.
\item Costello--Gwilliam 2017 \textit{Factorization Algebras in QFT}
  Vol~II \S 10--11: holomorphic factorisation algebras.
\item Francis 2013 arXiv:1211.5619 \S 2: $E_n$-chiral algebras and
  factorisation homology.
\end{itemize}

%--------------------------------------------------------------------
\subsection*{Cache appends (new patterns)}
%--------------------------------------------------------------------

Three pattern-level refinements of existing cache entries surface in
this assault; no new AP-CY entries are warranted, but the following
refinements belong at
\texttt{appendices/first\_principles\_cache.md} as AP-append rows:

\begin{enumerate}[leftmargin=1.5em]
\item \textbf{AP-CY144-refine (Wave-13 A1):} the eight
  Gritsenko--Cl\'ery diagonal-divisor forms correspond to eight
  $\mathrm{Sp}_{\Sigma_2,C}$-specialisations of $\Phih_3$, with
  $\Sigma_2=K3$-fibre at $N=1$ (canonical) and $\Sigma_2=H_N$
  (paramodular Humbert surface at level $N$) at
  $N\in\{2,3,4,6\}$, integer-weight sector. The three
  half-integer-weight specimens require extension of the
  specialisation functor beyond the integer Borcherds-lift tower;
  their specialisation cycles are not yet identified.
\item \textbf{AP-CY145-refine (Wave-13 A1):} the classical-limit
  BPS Lie subalgebra of $\mathrm{Sp}_{K3,E}(\Phih_3(K3\times E))$
  agrees with $\mathfrak g_{\Delta_5}$ at the denominator-identity
  level by direct Borcherds-lift computation
  (Gritsenko--Nikulin 1998), bypassing the Davison PBW route.
  Both routes produce Lie-algebra-level identification; the
  chiral-bialgebra-level upgrade
  $\mathbf H_{\Delta_5}\leftrightarrow\mathrm{Sp}_{K3,E}(\Phih_3)$
  remains open.
\item \textbf{AP-CY155 (new; Wave-13 A1, if append warranted):}
  ``Motivating-conjecture integrality boundary.'' The Lorgat 2020
  motivating conjecture extends a Borcherds-lift statement beyond
  the integer-weight paramodular tower; three of the eight
  Gritsenko--Cl\'ery specimens live at half-integer weight and
  require the Cl\'ery--Gritsenko 2008 Thm.~3.1 tower, whose
  BKM-correction is not constructed. The conjecture's scope must
  be qualified to $N\in\{1,2,3,4,6\}$ integer-weight regime, with
  the half-integer extension stated as a separate open problem.
\end{enumerate}

% End of Wave-13 A1 Gelfand rectification.
